\documentclass[11pt,a4paper]{article}
\usepackage[utf8]{inputenc}
\usepackage{amsmath}
\usepackage{amsfonts}
\usepackage{amssymb}
\usepackage{graphicx}
\usepackage{hyperref}
\usepackage{listings}
\usepackage{xcolor}

\definecolor{codegreen}{rgb}{0,0.6,0}
\definecolor{codegray}{rgb}{0.5,0.5,0.5}
\definecolor{codepurple}{rgb}{0.58,0,0.82}
\definecolor{backcolour}{rgb}{0.95,0.95,0.92}

\lstdefinestyle{mystyle}{
    backgroundcolor=\color{backcolour},   
    commentstyle=\color{codegreen},
    keywordstyle=\color{magenta},
    numberstyle=\tiny\color{codegray},
    stringstyle=\color{codepurple},
    basicstyle=\ttfamily\footnotesize,
    breakatwhitespace=false,         
    breaklines=true,                 
    captionpos=b,                    
    keepspaces=true,                 
    numbers=left,                    
    numbersep=5pt,                  
    showspaces=false,                
    showstringspaces=false,
    showtabs=false,                  
    tabsize=2
}

\lstset{style=mystyle}

\title{Q-Mini-WASM: A Hybrid Quantum-Classical Mixture-of-Experts Architecture for Deterministic In-Model Execution}
\author{}
\date{}

\begin{document}

\maketitle

\begin{abstract}
The historical trajectory of large language model (LLM) scaling has predominantly positioned transformer architectures as stochastic orchestrators of computation. In contemporary agentic workflows, when complex algorithmic reasoning or exact mathematical computation is required, models rely almost exclusively on external tool-use loops. This methodology involves generating a high-level code snippet, suspending autoregressive generation, transferring the payload to an isolated external sandbox or Python interpreter, and subsequently re-injecting the resultant output into the context window for further processing. While this paradigm has yielded state-of-the-art results on mathematical and coding reasoning evaluations, it imposes profound structural limitations. The external dependency introduces severe latency overheads, context fragmentation, and an inherent inability for the transformer to truly internalize the mechanical state transitions of the computation itself.
\end{abstract}

\section{Introduction}

The Q-Mini-WASM architecture represents a structural zenith in resolving this dichotomy. Developed as a highly specialized 500-million parameter prototype, this model operates at the absolute frontier of quantum machine learning (QML), discrete computational geometry, and low-level hardware orchestration. It explicitly internalizes exact computational execution by compiling a WebAssembly (WASM) interpreter directly into specific latent weights of a Sparse Mixture-of-Experts (MoE) framework. 

To achieve the zero-entropy state required by a machine instruction set architecture (ISA) without destroying the high-entropy semantic reasoning of natural language heads, Q-Mini-WASM deploys a highly heterogeneous hybrid topology. This comprehensive technical blueprint details the mathematical, quantum, and systems engineering foundations of the Q-Mini-WASM architecture. It formalizes the deployment of Parameterized Quantum Circuits (PQCs) to resolve discrete MoE routing, the application of modified Grover’s algorithms for unstructured ternary weight quantization, the computational geometry of Tropical Attention, the close-to-metal C++ SYCL hardware mapping required to execute this inference engine on Intel Alchemist and Battlemage graphics processing units, and the synthetic data pipeline required to bootstrap the model's fundamental understanding of computational execution.

\section{Pillar 1: The Quantum MoE Router (The Bridge)}

The fundamental challenge in constructing a hybrid transformer capable of executing deterministic WebAssembly instructions lies in the routing matrices. Classical MoE networks utilize dense feed-forward gating mechanisms terminating in softmax or top-k selection approximations to distribute tokens across available experts. While this continuous, differentiable methodology is sufficient for fuzzy, overlapping semantic representations, classical routing suffers from catastrophic failure modes when applied to rigid computational traces. Specifically, models succumb to routing collapse, where the network becomes heavily biased toward a small subset of experts, and token dropping, which occurs when expert computational capacities are strictly enforced and subsequently overloaded.

To enforce an execution-aware, hard-multiplexing control flow that routes discrete WASM tokens exclusively to dedicated execution experts while maintaining perfect load balancing, the Q-Mini-WASM architecture reformulates MoE routing as a discrete combinatorial optimization problem mapped to a quantum topology. By leveraging quantum tunneling and superposition, the architecture penetrates high-energy barriers in the loss landscape that would typically trap classical stochastic gradient descent (SGD) algorithms.

\subsection{Quadratic Unconstrained Binary Optimization (QUBO) Formulation}

Before quantum hardware can process the routing assignment, the optimization objective must be mathematically modeled as a Quadratic Unconstrained Binary Optimization (QUBO) problem. Let $T$ represent the set of input tokens in a micro-batch sequence, and $E$ represent the set of available experts within the MoE layer. The routing decision is parameterized as a discrete binary variable $x_{t,e} \in \{0, 1\}$, where $x_{t,e} = 1$ indicates that token $t$ is explicitly routed to expert $e$, and $x_{t,e} = 0$ indicates otherwise.

The primary objective is to maximize the cumulative token-to-expert affinity, represented as $\sum_{t,e} A_{t,e} x_{t,e}$, which is derived classically via a lightweight cross-attention dot product between the token contextual representation and the expert embedding vector. Because QUBO frameworks cannot natively enforce hard constraints through localized bounds, the strict routing requirements of the architecture must be seamlessly integrated into the objective function as quadratic penalty terms. Two fundamental constraints govern this network: the Top-K assignment constraint, dictating that every token must be routed to exactly $K$ experts, and the Expert Load Balancing constraint, ensuring no individual expert processes more than $C$ tokens per micro-batch to prevent buffer overflow.

Incorporating these constraints yields the comprehensive QUBO Hamiltonian objective:

$$H_Q = -\sum_{t,e} A_{t,e} x_{t,e} + \lambda_1 \sum_t \left( \sum_e x_{t,e} - K \right)^2 + \lambda_2 \sum_e \left( \sum_t x_{t,e} - C \right)^2$$

In this formulation, $\lambda_1$ and $\lambda_2$ are massive positive penalty scalars. Their magnitudes must be rigorously calibrated through hyperparameter tuning so that any deviation from the constraints incurs an energy penalty that thoroughly eclipses any possible optimization gain achieved by maximizing the raw affinity scores.

\subsection{Translating to the Ising Hamiltonian}

Quantum processing units (QPUs) fundamentally manipulate quantum states governed by Pauli operators. To execute the QUBO formulation on a quantum device, the binary decision variables $x_{t,e} \in \{0, 1\}$ must be algebraically mapped to quantum spin variables (the eigenvalues of the Pauli-Z operator), $s_i \in \{-1, 1\}$. This mapping is achieved through the linear transformation:

$$x_i = \frac{1 - s_i}{2}$$

Substituting this transformation into the expanded QUBO equation and isolating the linear and quadratic terms translates the classical combinatorial problem into the Cost Hamiltonian $H_C$, representing an Ising spin-glass model:

$$H_C = \sum_{i<j} J_{i,j} Z_i Z_j + \sum_i h_i Z_i$$

Here, $Z_i$ and $Z_j$ denote the Pauli-Z operators acting on the specific qubits representing the flattened coordinate routing variables. The interaction strength term $J_{i,j}$ encapsulates the physical entanglement required to enforce the capacity and top-k penalties, derived directly from the cross-terms of the squared QUBO penalty functions. The parameter $h_i$ represents the local magnetic field applied to each qubit, factoring in both the classical token affinities $A_{t,e}$ and the linear components of the expanded penalty functions. The lowest eigenvalue (the ground state) of this Hamiltonian mathematically encodes the optimal, perfectly load-balanced MoE routing matrix.

\subsection{Angle Embedding and Barren Plateau Mitigation}

To extract the ground state of $H_C$, the Quantum Approximate Optimization Algorithm (QAOA) is deployed within the hybrid loop. QAOA constructs a Parameterized Quantum Circuit (PQC) based on adiabatic quantum computation, applying a sequence of alternating unitary operators: one derived from the Cost Hamiltonian $e^{-i\gamma H_C}$, and one from a non-commuting Mixer Hamiltonian $e^{-i\beta H_B}$, where $H_B = \sum_i X_i$ represents the Pauli-X operator. The role of the Mixer Hamiltonian is to induce quantum superposition, allowing the state to transition between different routing topologies through quantum tunneling, preventing the system from freezing in a suboptimal local minimum.

However, the physical constraints of the Noisy Intermediate-Scale Quantum (NISQ) era strictly prohibit the loading of a 4096-dimensional continuous latent space directly into quantum memory using basic basis encoding. Attempting to map 4096 features directly to individual qubits would require a lattice structure geometrically impossible on current 133-qubit to 156-qubit superconducting processors.

To bridge this dimensional gap, the architecture employs a hierarchical dimensionality reduction protocol combined with Dense Angle Embedding. A classical dense neural network first projects the 4096-dimensional token state into a highly compressed 8-dimensional manifold. These 8 dimensions are subsequently encoded into an 8-qubit quantum register. To maximize quantum phase space utility and achieve exponential data compression, the classical features are mapped directly to the rotation parameters (e.g., $R_y(\theta)$) of the individual qubits.

Furthermore, deep PQCs—such as the QAOA circuit required for intricate MoE routing—are exceptionally vulnerable to the barren plateau phenomenon. Driven by the concentration of measure and expressibility defined by Weingarten calculus, the variance of the cost function's gradient with respect to any circuit parameter $\theta_k$ vanishes exponentially as the number of qubits ($N$) or the depth of the circuit increases:

$$Var[\partial_{\theta_k} C] \in \mathcal{O}\left(\frac{1}{2^N}\right)$$

In a barren plateau, the optimization landscape becomes catastrophically flat. Classical optimizers like AdamW receive gradient signals that are indistinguishable from statistical hardware noise, completely halting the training of the quantum router. To enforce convergence during hybrid training, three advanced mitigation strategies are explicitly layered into the Q-Mini-WASM framework:

\begin{itemize}
    \item \textbf{Local Cost Functions:} Measures localized qubit observables that only evaluate adjacent correlations rather than global states across the heavy-hex lattice. Transitions the gradient scaling from exponentially vanishing to polynomially vanishing, preserving signal integrity.
    \item \textbf{Layer-wise Pretraining:} Initializes shallow, one-block circuits trained to localized convergence before incrementally appending additional parameterized unitary layers. Prevents the deep entanglement structure from overwhelming initial gradient signals; avoids starting in a barren plateau.
    \item \textbf{Lie Algebraic Subspaces:} Geometrically constrains the dynamical Lie algebra of the QAOA Cost Hamiltonian to specific algebraic subspaces. Naturally maintains non-zero gradient expectations regardless of circuit depth by preventing uniform mixing in the full Hilbert space.
\end{itemize}

\subsection{PyTorch and PennyLane Hybrid Implementation}

The following conceptual implementation utilizes the PennyLane framework integrated natively with the PyTorch autograd engine. During the backward pass, gradients are passed from the QPU back to the classical GPU utilizing the parameter-shift rule, extracting exact analytical derivatives without suffering from finite-difference shot-noise degradation.

\begin{lstlisting}[language=Python]
import torch
import torch.nn as nn
import pennylane as qml
from pennylane import numpy as np

NUM_QUBITS = 8
dev = qml.device('default.qubit', wires=NUM_QUBITS)

def qaoa_layer(gamma, beta, compressed_affinities, penalty_factor):
    """Applies a single adiabatic block of the QAOA unitary evolution."""
    # Execute Cost Hamiltonian (Ising Model H_C)
    for i in range(NUM_QUBITS):
        # Apply local magnetic fields encoding classical semantic projections
        qml.RZ(gamma * compressed_affinities[i], wires=i)
        
    # Execute Ising J_ij interactions to enforce capacity and top-k penalties
    for i in range(NUM_QUBITS):
        for j in range(i + 1, NUM_QUBITS):
            qml.IsingZZ(gamma * penalty_factor, wires=[i, j])
            
    # Execute Mixer Hamiltonian H_B (Pauli-X Rotations for quantum tunneling)
    for i in range(NUM_QUBITS):
        qml.RX(beta, wires=i)

@qml.qnode(dev, interface='torch', diff_method='parameter-shift')
def quantum_router_circuit(gammas, betas, compressed_affinities, layers=3):
    """Constructs the QAOA PQC to evaluate mathematically optimal routing matrices."""
    # Initialization: Prepare equal superposition state |+>
    for i in range(NUM_QUBITS):
        qml.Hadamard(wires=i)
        
    # Dense Angle Embedding for Dimensionality Reduction
    for i in range(NUM_QUBITS):
        qml.RY(compressed_affinities[i], wires=i)
        
    # Adiabatic Evolution: Apply alternating unitary blocks
    for p in range(layers):
        qaoa_layer(gammas[p], betas[p], compressed_affinities, penalty_factor=2.5)
        
    # Measurement: Extract expectation values of Pauli-Z (Spin states)
    return [qml.expval(qml.PauliZ(i)) for i in range(NUM_QUBITS)]

class HybridQuantumMoE(nn.Module):
    def __init__(self, d_model=4096, num_experts=8, qaoa_layers=3):
        super(HybridQuantumMoE, self).__init__()
        # Classical projection network compressing 4096 dimensions to 8 qubits
        self.compressor = nn.Linear(d_model, NUM_QUBITS)
        
        # Trainable Quantum Angles initialized via identity block to mitigate plateaus
        self.gammas = nn.Parameter(torch.randn(qaoa_layers) * 0.01)
        self.betas = nn.Parameter(torch.randn(qaoa_layers) * 0.01)
        
        self.experts = nn.ModuleList([nn.Linear(d_model, d_model) for _ in range(num_experts)])

    def forward(self, hidden_states):
        # Phase 1: Compress high-dimensional state for NISQ encoding
        compressed_state = torch.tanh(self.compressor(hidden_states)).squeeze()
        
        # Phase 2: QPU Execution. QAOA circuit computes discrete combinatorial routing
        q_routing_exp = quantum_router_circuit(self.gammas, self.betas, compressed_state)
        
        # Phase 3: Map Pauli-Z expectation values [-1, 1] to binary routing probabilities
        q_routing_probs = [(val + 1.0) / 2.0 for val in q_routing_exp]
        routing_weights = torch.stack(q_routing_probs)
        
        # Phase 4: Sparse Expert Execution
        out = torch.zeros_like(hidden_states)
        for i, expert in enumerate(self.experts):
            out += routing_weights[i] * expert(hidden_states)
            
        return out
\end{lstlisting}

\section{Pillar 2: Ternary Quantization \& Grover's Search}

While probabilistic semantic routing seamlessly directs natural language sequences toward dense Multi-Layer Perceptrons (MLPs) operating at standard FP16 precision, the dedicated WASM execution experts within Q-Mini-WASM require an environment of absolute, uncompromising mathematical rigidity. To guarantee deterministic execution and perfectly simulate exact processor stack mechanics over million-step traces, the architecture forces the underlying parameter weights of these specific WASM execution experts into an unstructured ternary state: $W \in \{-1, 0, 1\}$.

Standard linear Min-Max quantization algorithms—routinely utilized by frameworks like Unsloth—induce a fatal phenomenon known as floating-point drift. In a semantic network, minor bit-crushing noise merely results in the occasional selection of a valid linguistic synonym. In WASM state emulation, shifting a parameter by even $1e-4$ will geometrically distort the topological boundary of a logical gate, causing the neural network to fetch the incorrect physical integer from the emulated stack, thereby crashing the deterministic execution trace instantly.

The WASM expert weights must be explicitly trained and retained in exact discrete states.

\subsection{The Straight-Through Estimator (STE) in Classical Forward Passes}

Classical neural network training mechanisms, predicated on backpropagation, are violently disrupted by ternary quantization. Backpropagation of error gradients requires continuous, smooth, differentiable manifolds; attempting to backpropagate through hard ternary step functions or clamping mechanisms yields analytical gradients of exactly zero almost everywhere, completely trapping the classical optimization trajectory in stalled flat regions.

To navigate this limitation during the classical phases of the hybrid training curriculum, Q-Mini-WASM employs the Straight-Through Estimator (STE). The STE bypasses the non-differentiable quantization step during the backward pass. It estimates the gradients by treating the non-differentiable rounding or thresholding function as an identity operator, passing the incoming gradient directly back to the latent continuous weights uncorrupted.

The structural implementation in PyTorch requires a custom forward function that carefully manipulates the computational graph utilizing the detach() method. This trick cleanly isolates the discrete mathematical calculation from the continuous derivative tracking required by the autograd engine.

\begin{lstlisting}[language=Python]
import torch
import torch.nn as nn
import torch.nn.functional as F

class TernaryWASMExpert(nn.Linear):
    def __init__(self, in_features, out_features, bias=False, precision="ternary"):
        super(TernaryWASMExpert, self).__init__(in_features, out_features, bias)
        self.precision = precision
        # Sustain the variance during initialization for ternary parameters
        bound = (2.0 / (0.75 * in_features)) ** 0.5 
        nn.init.uniform_(self.weight, a=-bound, b=bound)

    def binarize_ternary(self, weight):
        """Forces continuous latent weights into strict {-1, 0, 1} states."""
        # Calculate the absolute mean for adaptive thresholding
        abs_mean = weight.abs().mean()
        
        # Discretize the weights
        weight_ternary = torch.round(weight / abs_mean)
        weight_ternary = torch.clamp(weight_ternary, -1.0, 1.0)
        
        # The Straight-Through Estimator (STE) Graph Manipulation
        # During the forward pass, the actual discrete `weight_ternary` is utilized.
        # During the backward pass, the gradient routes around `weight_ternary` 
        # and flows directly into the continuous `weight`.
        return (weight_ternary - weight).detach() + weight

    def forward(self, input_tensor):
        if self.precision == "ternary":
            discrete_weight = self.binarize_ternary(self.weight)
        else:
            discrete_weight = self.weight
            
        return F.linear(input_tensor, discrete_weight, self.bias)
\end{lstlisting}

\subsection{Grover's Search and the Register Counting (RC) Oracle}

While the STE algorithm successfully allows classical gradient descent to proceed without zeroing out, it fundamentally disconnects the true discrete weight mapping from the gradient trajectory, leading to a profound mathematical disconnect that frequently drives the network into shallow local minima. To derive the globally optimal ternary parameter configuration without relying on continuous probability approximations, Q-Mini-WASM reframes the ternary weight optimization as a vast unstructured search problem through a discrete combinatorial space. By utilizing quantum superposition via a modified Grover’s Search algorithm, the pipeline locates the optimal parameter configuration with a quadratic speedup $O(\sqrt{N})$ over classical brute-force evaluation $O(N)$.

Because standard quantum hardware utilizes binary qubits (base-2), mapping a base-3 ternary weight structure necessitates a specialized dual-qubit encoding protocol per individual network parameter. An individual ternary weight $w$ is mapped onto a two-qubit state $|q_1 q_0\rangle$ utilizing the following logic:
\begin{itemize}
    \item $w = 0 \rightarrow |00\rangle$
    \item $w = 1 \rightarrow |01\rangle$
    \item $w = -1 \rightarrow |10\rangle$
    \item $|11\rangle \rightarrow \text{Invalid State}$
\end{itemize}

The Grover initialization phase prepares an equal superposition of all geometrically valid ternary configurations by applying specifically tailored Hadamard and controlled phase operators designed to suppress the generation of the invalid $|11\rangle$ state across the register, establishing the initial search state $|s\rangle$.

The functional core of Grover's Algorithm relies on a Quantum Oracle operator $U_\omega$, an intricate quantum arithmetic circuit designed to evaluate the validity of a given superimposed state and mark it via a precise phase inversion. For the Q-Mini-WASM execution expert, this Oracle must execute a full neural network forward pass entirely in quantum superposition, compute the cross-entropy loss against a classically provided training batch, and mark the ternary configurations that produce a loss falling below a dynamically specified threshold.

Classical approaches to quantum algorithm design frequently rely on the Quantum Phase Estimation (QPE) Oracle. However, QPE requires repeated implementations of the full training dataset for each precision qubit, resulting in an exponentially scaling computational depth that inevitably induces total decoherence on NISQ-era hardware.

Consequently, the Q-Mini-WASM architecture mandates the construction of a Register Counting (RC) Oracle. The RC Oracle utilizes a quantum incrementing adder circuit to maintain a direct binary count of correct training predictions. Its circuit depth scales strictly linearly with the dataset size, rendering simulation and physical execution highly viable.

Within the RC Oracle framework, an ancillary quantum register $|c\rangle$ is initialized to $|0\rangle$. The circuit processes the training batch sequentially. For every sample $(x_i, y_i)$, a unitary representation of the dense neural network $U_{NN}$ is applied to the superposed weight register $|W\rangle$ and the input data $|x_i\rangle$, outputting a superposed prediction state. A controlled incrementer circuit ($C-INC$) then evaluates the prediction state against the embedded ground truth state $|y_i\rangle$. If the prediction precisely matches the ground truth, the amplitude state of the counting register $|c\rangle$ is augmented by $+1$.

Following iteration through all sequence samples, a final quantum comparator arithmetic block evaluates the condition $|c\rangle > \text{threshold}$ (where threshold is the predefined cross-entropy equivalent). If the condition is met, a multi-controlled Pauli-Z gate applies a negative phase shift strictly to the valid state amplitudes.

Following the Oracle marking, the Grover Diffusion Operator (an inversion about the mean, represented as $U_s = 2|s\rangle\langle s| - I$) is applied. This operator geometrically amplifies the probabilities of the phase-inverted optimal ternary states while destructively interfering with sub-optimal configurations. After $\approx \frac{\pi}{4}\sqrt{N}$ iterations, measurement of the register collapses the quantum system precisely onto the optimal ternary weight matrix that minimizes the expert sub-network loss, entirely circumventing the vanishing gradients of classical backpropagation.

\section{Pillar 3: The Geometric HullKVCache}

The context window capability of standard LLMs—such as the 200,000-token capacity of the original MiniMax M2.5 architecture—relies extensively on Lightning Attention, a linear attention mechanism that resolves the quadratic complexity bottleneck of standard softmax transformers. By stripping away the non-linear softmax operation in favor of recurrent linear right-multiplication ($(Q \cdot K^T) \cdot V \rightarrow Q \cdot (K^T \cdot V)$), linear attention achieves length-invariant $O(1)$ computational complexity.

However, linear attention inherently relies on iteratively updating an accumulated, exponentially decayed state matrix ($S_t = \lambda S_{t-1} + K_t^T V_t$) to maintain causality. While this exponential decay is exceptionally memory-efficient and ideal for probabilistic natural language retrieval, it is fundamentally destructive to the exact, discrete machine states required by a WebAssembly execution trace. A mathematically decayed and smoothed state matrix cannot exactingly recover a specific machine integer pushed to a WASM stack 50,000 tokens prior.

\subsection{Computational Geometry of 2D Attention Heads}

To resolve this paradox without reverting to quadratic bottlenecks, the Q-Mini-WASM architecture executes a surgical modification. A specific subset of the attention heads are carved out to serve as the WASM execution heads ($H_{WASM}$). These heads are forced into an extreme dimensionality constraint: they are restricted to a strictly two-dimensional subspace ($d_{head} = 2$).

This dimensionality reduction is not a lossy compression technique; rather, it is a foundational geometric transformation that lifts the attention kernel into tropical projective space. In this $\mathbb{TP}^2$ space, the attention mechanism is modeled using Tropical Geometry and max-plus algebra. Within the max-plus semiring, the probabilistic dense matrix multiplication defined by the standard Euclidean softmax function is abandoned. Instead, the standard addition is replaced by the maximum operator ($\oplus = \max$), and standard multiplication is replaced by addition ($\otimes = +$).

Consequently, the attention operation transitions into a deterministic geometric search for the boundary of a convex polygon. The Multi-Head Tropical Attention (MHTA) acts as a universal approximator of dynamic programming and tropical transitive closure. By evaluating reasoning piecewise-linearly and maintaining 1-Lipschitz properties, this architecture ensures that the induced polyhedral decision boundaries remain razor-sharp and scale-invariant, completely negating the smoothing artifacts generated by classical linear decay.

Rather than computing a continuous weighted sum over all historical tokens, the model queries the upper convex hull of all historical Key vectors. Geometrically, the exact output deterministic state is retrieved by locating the extreme boundary point on the upper convex hull in the precise direction defined by the incoming Query vector.

\subsection{The Monotone-Chain Update Algorithm and $O(\log N)$ Binary Search}

This geometric transformation necessitates a specialized inference data structure: the HullKVCache. Because finding a tangent on a dynamically maintained convex hull operates in logarithmic time, the HullKVCache enables an $O(\log N)$ convex-hull query. This allows the execution of deterministic WASM traces at generation speeds exceeding 30,000 tokens per second, entirely circumventing the $O(N^2)$ attention bottleneck while preserving absolute precision.

As the WASM execution trace progressively generates new tokens, the system continuously updates the geometric boundary utilizing a streaming variant of the Monotone-Chain algorithm (Andrew's algorithm). When a new execution token projects a Key vector $P_k$ into the 2D plane, the custom Triton kernel evaluates if the new coordinate geometrically engulfs previous points on the hull.

The kernel calculates the cross product of the vectors formed by the last two historical points ($P_{k-2}, P_{k-1}$) and the new point ($P_k$):

$$\text{Cross} = (x_{k-1} - x_{k-2})(y_k - y_{k-2}) - (y_{k-1} - y_{k-2})(x_k - x_{k-2})$$

If the cross product returns a positive value (or zero), indicating a non-convex inward orientation, the historical point $P_{k-1}$ is aggressively popped from the upper hull boundary, as it no longer forms the strict convex outer shell.

Once the convex hull is synchronized via barrier operations, the system must execute the max-plus projection. Because the vertices of the monotone-chain hull are strictly sorted along the boundary, the hardware threads execute a highly parallel binary search in $O(\log H)$ time. The search unrolls a bounded loop to calculate the local slope derivative between adjacent mid-points on the hull, quickly dictating if the optimal tangent lies to the left or the right. Upon locating the extreme tangent vertex that maximizes the max-plus projection against the Query vector $Q$, the kernel retrieves its corresponding high-dimensional Value vector $V$ to populate the residual stream, perfectly recovering the uncorrupted historical machine state.

\subsection{Convex Hull Trace Eviction and Infinite Contexts}

Storing sequence states for an unbroken, multi-million step algorithm execution trace presents a catastrophic physical memory limitation. However, the unique geometric properties of the HullKVCache offer a mathematically guaranteed solution to this bottleneck: Convex Hull Trace Eviction.

In the mathematics of max-plus algebra, the query output is dictated exclusively by the extreme boundary points. Any historical Key vector $P_i$ that falls strictly into the interior geometry of a newly updated hull becomes mathematically irrelevant to all future queries; an internalized point can never represent the maximum extreme point in any subsequent projection.

Consequently, the memory footprint of the context window is decoupled from the raw linear sequence length $N$ and becomes entirely bounded by $O(H)$, where $H$ is the number of vertices actively forming the upper convex hull. Because the expected number of vertices on a random convex hull scales logarithmically, a massive 5-million token execution trace is geometrically compressed down to fewer than 5,000 active boundary tokens. The high-dimensional $V$ vectors of internalized points are aggressively evicted from High Bandwidth Memory (HBM), allowing the context window to theoretically scale infinitely without triggering an Out-Of-Memory (OOM) exception.

\section{Pillar 4: Intel ARC / SYCL Hardware Mapping}

Deploying a highly heterogeneous architecture—which must seamlessly interleave continuous probabilistic MoE routing, rigid unstructured ternary quantum matrices, and deterministic 2D geometric convex-hull searches—demands execution close-to-metal, bypassing standard high-level AI frameworks. The primary deployment targets for the Q-Mini-WASM inference engine are Intel ARC graphics processing units, specifically the first-generation Alchemist (Xe-HPG) and second-generation Battlemage (Xe2-HPG) microarchitectures.

The engine leverages the Intel oneAPI Data Parallel C++ (DPC++/SYCL) programming model to explicitly map architectural abstractions to the physical silicon execution units.

\subsection{SYCL Execution Hierarchy and Xe-Core Configurations}

The Intel Xe GPU family consists of hierarchical building blocks designed for both dense matrix computations and highly parallelized vector operations. To achieve peak device occupancy, the SYCL execution hierarchy is meticulously mapped to the hardware:
\begin{enumerate}
    \item \textbf{Work-items:} Represent a single thread of execution, mapped directly to an individual SIMD lane.
    \item \textbf{Sub-groups:} A collection of work-items executing in lockstep, mapping directly to a Vector Engine hardware thread.
    \item \textbf{Work-groups:} A representation of related work-items executing exclusively on a single Xe-core.
\end{enumerate}

The transition from the Alchemist to the Battlemage microarchitecture fundamentally dictates the thread mapping strategy. Alchemist features a dual subslice design comprising 16 Xe Vector Engines (XVE, SIMD8 256-bit width) and 16 Xe Matrix Extensions (XMX, 1024-bit width) per core. Battlemage abandons this for 8 Vector Engines and 8 Matrix Engines per core. Crucially, Battlemage doubles the Arithmetic Logic Units (ALUs) to a native SIMD16 (512-bit width), vastly increasing single-thread instruction-level parallelism and fine-grained branching execution, while upgrading the XMX engines to massive 2048-bit wide arrays to prevent systolic starvation.

\begin{table}[h]
\centering
\begin{tabular}{|l|l|l|p{5cm}|}
\hline
\textbf{Architectural Parameter} & \textbf{Arc Alchemist (Xe-HPG)} & \textbf{Arc Battlemage (Xe2-HPG)} & \textbf{Inference Engineering Implication} \\
\hline
Vector Engines per Core & 16 (SIMD8, 256-bit) & 8 (SIMD16, 512-bit) & SYCL sub-group mapping must adapt dynamically; Battlemage handles complex branching operations more efficiently. \\
\hline
XMX Matrix Engines & 16 (1024-bit width) & 8 (2048-bit width) & DPAS payload packing strategies must be adjusted to feed wider arrays without causing execution stalls. \\
\hline
Shared Local Memory (SLM) & 128 KB per Xe-core & 192 KB per Xe-core & Battlemage enables larger continuous HullKVCache convex boundaries to be pinned directly to the L1\$. \\
\hline
\end{tabular}
\end{table}

\subsection{Asynchronous Dispatch Strategy: XMX vs. Vector Engines}

To optimize MoE execution, the Q-Mini-WASM engine exploits a profound physical advantage of the Intel Xe architecture: the capability to permit concurrent execution of Vector Engines (XVE) and Matrix Engines (XMX) within the identical Xe-core boundary.

\textbf{Vector Engines (XVE) for Logic-Heavy Routing:} The sparse gating network—which calculates classical token affinities and executes the execution-aware hard-switch for WASM tokens—is highly branch-dependent and bound by memory bandwidth. These operations are dispatched exclusively to the Vector Engines utilizing the Explicit SIMD SYCL Extension (ESIMD). ESIMD enables close-to-metal programming by bypassing unpredictable compiler heuristics and standard SYCL implicit vectorization. By defining a logical sub-group size of exactly 1, the kernel maps directly to a single hardware thread, allowing developers to manually author explicit cross-lane data sharing code.

\textbf{Matrix Engines (XMX) for MLP Execution:} Once the routing indices are determined by the ESIMD kernels, the compute-bound dense matrix multiplications of the selected expert blocks are instantly dispatched to the XMX systolic arrays. This completely bypasses heavy libraries like oneDNN in favor of the specialized SYCL \texttt{sycl\_ext\_oneapi\_matrix} (joint matrix) extension. This extension natively targets the Intel hardware's Dot Product and Accumulate Systolic (DPAS) instructions.

The DPAS workflow initiates 2D matrix representations directly within registers shared among a sub-group. The \texttt{joint\_matrix\_load} function streams the weights and continuous token representations from HBM, \texttt{joint\_matrix\_mad} executes the fused multiply-add operation across the interconnected grid of processing elements, and \texttt{joint\_matrix\_store} writes the accumulated outputs back to global memory.

\subsection{C for Metal (CM) and Ternary Weight Packing}

The 8 dedicated WASM execution experts are restricted entirely to the discrete ternary space $W \in \{-1, 0, 1\}$. A single ternary digit ("trit") holds exactly $\log_2(3) \approx 1.585$ bits of information. Storing a single trit inside a standard 8-bit INT8 format wastes vast amounts of VRAM bandwidth. Therefore, the architecture employs optimal mathematical block packing. A block size of 5 trits yields $3^5 = 243$ distinct states. Because 243 is comfortably close to $2^8 = 256$, exactly 5 ternary weights can be packed into a single 8-bit byte, achieving an information density operating at 99.06\% of the theoretical perfection limit.

However, the XMX systolic arrays require continuous streams of natively formatted data types (FP16/INT8) to perform DPAS operations. The 8-bit payloads must be decompressed back into 5 independent data structures dynamically on the fly. Standard SYCL abstractions generate fatal latency penalties during this unpacking phase.

To execute bitwise unpacking microkernels at maximum throughput, the engine employs Intel C for Metal (CM), a specialized GPU programming language designed as an assembler equivalent. CM provides an explicit SIMD model that bypasses standard C++ overheads in favor of direct hardware intrinsics. Operating in explicit SIMD, the CM kernels use the \texttt{rdregion} (read region) intrinsic to extract sub-matrices and pull packed bytes into the execution lanes, simulating integer division via advanced bitwise shifting and masking logic. Once decompressed, the values are written back using the \texttt{wrregion} intrinsic.

To prevent register spilling to scratch memory—which incurs devastating latency penalties—the Intel DPC++ compiler (icpx) is forced into large GRF mode via the \texttt{-ftarget-register-alloc-mode=large} flag. This grants a single hardware thread access to 256 General Register File (GRF) registers instead of the default 128, easily absorbing the intense register pressure generated by the CM ternary unpacking sequence and allowing the data to feed the XMX arrays directly.

\subsection{Driver-Level Memory Paging for 5M Token Contexts}

Executing complex self-contained deterministic algorithms natively requires execution traces that scale up to 5 million tokens. Even with the geometric $O(H)$ compression provided by Convex Hull Trace Eviction, storing the continuous sequence states requires robust backend memory management.

The Q-Mini-WASM engine leverages the bleeding-edge memory capabilities of the Intel Xe Linux Kernel Driver (i915/xe). By integrating directly with the driver's Shared Virtual Memory (SVM) subsystem, the host CPU and GPU share a unified virtual address space. When the SYCL HullKVCache kernel identifies that a block of historical tokens has been permanently subsumed into the interior of the geometric polygon, the user-mode driver interfaces with the kernel-mode driver's garbage collector. Using the \texttt{drm\_gpusvm\_devmem} routines, the GPU seamlessly unmaps and destroys the page table bindings for those specific evicted ranges, allowing the high-dimensional Value vectors to be aggressively paged back to the host system's DRAM.

To eliminate the CPU latency and execution stalling incurred during these massive page fault handling events, the engine utilizes the driver's Transparent Hugepages (THP) support. Utilizing the \texttt{MIGRATE\_VMA\_SELECT\_COMPOUND} flag, the driver operates at the macro-folio level, migrating memory via 2MB hugepages rather than looping over individual 4KB blocks. This architectural pairing accelerates SVM page fault servicing by a factor of 7—dropping the total latency to an imperceptible 132 microseconds—ensuring the engine can unroll 5-million-step WebAssembly traces natively without ever starving the execution pipelines.

\section{Pillar 5: The Synthetic Data Pipeline}

Transforming the Q-Mini-WASM framework from a theoretical topological construct into a functional inference engine requires an exhaustive pre-training curriculum. The transformer must physically internalize the rigid mechanics of the WebAssembly instruction set architecture (ISA), explicitly recognizing the step-by-step state modifications of programmatic memory structures.

\subsection{Wasmtime Instrumentation and Stack/Memory Deltas}

The bootstrapping training corpus must map high-level algorithmic logic (C, C++, Rust) to its compiled WASM equivalent, and subsequently flatten that logic into strictly stepped token execution traces. Generating this dataset at scale requires bridging deterministic classical algorithms—such as exact Hungarian algorithm solvers, Dijkstra's shortest path routines, or Arto Inkala Sudoku resolvers—into the synthetic generation pipeline.

This is accomplished by leveraging the Wasmtime runtime engine, manipulated via advanced Python API bindings. Wasmtime provides deep, low-level introspection into the WASM execution sandbox. A custom Python tracer utilizes these bindings to hook directly into the WebAssembly linear memory. WebAssembly linear memory exists as a contiguous array of bytes (allocated in multiples of a 64-kilobyte page size) representing global data, statics, and shadow stacks.

The instrumentation pipeline compiles the C-algorithms into WASM intermediate representations using LLVM. By enforcing epoch-based interruptions rather than fuel-based metering, the tracer achieves fine-grained deterministic interception without incurring unacceptable latency overheads. After the execution of every individual WASM opcode (e.g., \texttt{i32.add}, \texttt{i32.store}, \texttt{br\_if}), the instrumentation hooks extract the exact differential delta of the linear memory arrays and the execution stack.

This multidimensional machine state is then flattened into strict textual sequences comprising augmented WASM-specific operational and state-management tokens that the transformer can ingest.

\begin{table}[h]
\centering
\begin{tabular}{|l|l|l|p{4cm}|}
\hline
\textbf{Algorithmic Domain} & \textbf{High-Level Code} & \textbf{WASM Intermediate} & \textbf{Flattened Token Format} \\
\hline
Variable Initialization & \texttt{int counter = 5;} & \texttt{i32.const 5} & \texttt{05 00 00 00 commit(+1,sts=1,bt=0)} \\
\hline
Stack Arithmetic & \texttt{counter + 3;} & \texttt{i32.const 3; i32.add} & \texttt{03 00 00 00 commit(+1); 08 00 00 00 commit(-1,sts=1)} \\
\hline
Memory Mutation & \texttt{*ptr = val;} & \texttt{i32.store ptr\_adr val\_adr} & \texttt{mem\_write(sts=1,bt=0)} \\
\hline
Conditional Control Flow & \texttt{if (counter > 0)} & \texttt{i32.gt\_s; br\_if cmp\_res\_1} & \texttt{br\_target(jump\_if\_true, offset=12)} \\
\hline
\end{tabular}
\end{table}

\subsection{Intentional Fault Injection and State-Recovery}

An LLM trained strictly on perfect, untainted execution traces remains exceptionally brittle. A single hallucinated integer or a slight miscalculation generated during a multi-million step autoregressive rollout will lead to cascading, unrecoverable trace failures, permanently crashing the emulated program. To build mechanical resilience, the pipeline employs intentional Fault Injection strategies.

During the Wasmtime trace generation, the Python instrumentation intentionally introduces localized corruption events. These faults simulate the bounds of adversarial perturbation: flipping bits within a newly written 32-bit integer, intentionally dropping an item from the emulated stack, or simulating an out-of-bounds pointer calculation.

The generated token trace is subsequently appended with specialized error-handling routines and state-recovery token sequences. When encountering these corrupted traces during training, the 2D execution heads are forced to recognize the state discrepancy. Utilizing the $O(\log N)$ convex-hull query of the HullKVCache, the model learns to geometrically retrieve the last known valid state of the linear memory stored deeper in the context window. The model must then generate the explicit corrective tokens required to revert the pointer arithmetic or fix the stack alignment, thereby internalizing state-recovery mechanisms directly into the parameter weights.

\subsection{Continuous Pre-training and CISPO Integration}

The final integration of this synthetic data relies on a rigorous, multi-phase Continuous Pre-training (CPT) and Reinforcement Learning curriculum designed to protect the model's existing semantic knowledge base.

\begin{enumerate}
    \item \textbf{Latent ISA Mapping:} The base M2.5 semantic experts (0-255) are strictly frozen. A high-rank Low-Rank Adaptation (LoRA) module is applied to the semantic Lightning Attention matrices to bridge the natural language prompt with the initiation of WASM logic. Concurrently, the 2D Tropical Attention heads and the 8 dedicated ternary WASM experts are fully unfrozen and trained aggressively via teacher forcing to memorize the physical mechanics of stack pushing, popping, and binary arithmetic.
    \item \textbf{Autoregressive Trace Unrolling:} The model engages in open-ended trace unrolling without teacher forcing, ensuring the 2D convex-hull queries correctly and reliably retrieve distant historical stack states across the sequence.
    \item \textbf{CISPO Reinforcement Learning:} The final phase applies rigorous alignment utilizing the Clipped IS-weight Policy Optimization (CISPO) reinforcement learning algorithm. Standard Proximal Policy Optimization (PPO) notoriously destabilizes MoE routing matrices over extended traces by destroying gradient flow to under-utilized experts. CISPO mitigates gradient norm explosions by clipping the importance sampling (IS) weights instead of the policy updates. This maintains continuous, stable gradient contributions from all tokens across the execution trace, ensuring the quantum-derived routing topologies and discrete ternary weights remain flawlessly calibrated.
\end{enumerate}

\section{Synthesized Architectural Conclusions}

The Q-Mini-WASM architecture marks a decisive departure from the stochastic paradigms governing contemporary large language models. By explicitly isolating high-entropy semantic reasoning from the zero-entropy rigidity of deterministic computation, it internalizes the WebAssembly instruction set architecture natively within its own neural weights.

This synthesis relies on the seamless convergence of disparate computational disciplines. The reformulation of MoE routing as an Ising Hamiltonian solved via Parameterized Quantum Circuits (PQCs) and the derivation of discrete ternary experts through modified Grover's algorithms establish a mathematically pristine logic framework immune to continuous floating-point drift. Concurrently, the geometric transformation of the attention mechanism into a two-dimensional tropical projective space replaces classical linear decay with $O(\log N)$ max-plus convex-hull queries. This ensures that million-step algorithm traces are executed with absolute precision while maintaining an infinitely scalable context window constrained only by active hull boundaries. By aggressively orchestrating these structures at the lowest hardware levels—utilizing SYCL joint matrix extensions, C for Metal intrinsics, and Intel Xe driver paging mechanisms—Q-Mini-WASM provides the definitive blueprint for the generation of true Large Reasoning Models capable of functioning as their own deterministic operating environments.

\end{document}